% !TeX root = ../教材题库_全部条目_转写工作稿.tex
% Source: XJ-B1-C04-EX-P120-Y0749
% 第4章 / 4.3 对数函数 / 教材 p.120 / PDF 第127页
% OCR未校对：公式与文字均可能有误；未完成数学TeX转写；禁止直接发给学生
% 原题及相邻原解截图（核对用）：../题目原图/XJ-B1-C04-EX-P120-Y0749-01.webp
% 本文件只在转写完成后作为题干使用；OCR 条目尚不是可靠题干。
\begin{bankitem}{4.3 例题 例题9}{第4章 / 4.3 对数函数 / 教材 p.120 / PDF 第127页}{XJ-B1-C04-EX-P120-Y0749}{OCR未校对：公式与文字均可能有误；未完成数学TeX转写；禁止直接发给学生}
\begin{ocrdraft}
[教材p.120 / PDF第127页]
\par
0二0]地震的强烈程度通常用里氏震级M=lgA—lgA0表示，这里A是距离震
\par
中100km处所测量地震的最大振幅，AO是该处的标准地震振幅．
\par
（D若一次地震测得A=25mm，一0．佣1mm，则该地震的震级是多少
\par
()g2．50，4）？
\par
（2）计算里氏8级地震的最大振幅是里氏5级地震最大振幅的多少倍．
\par
解（1）M一lg25\ocrglyph{00AF}lg0．001
\par
0，001
\par
一lg25000
\par
一lg2，5十lg101
\par
因此，该地震的震级约为里氏4．4级．
\par
里氏震级表
\par
9
\par
（2）设里氏8级和5级地震的最大振幅分别为AI，A2．
\par
由题意，得
\par
AI
\par
AI
\par
A2
\par
A2
\par
lg
\par
7
\par
6
\par
甚强
\par
由上可得，
\par
4
\par
0
\par
极微
\par
=103=1000．
\par
\ensuremath{\cdot}105
\par
因此，里氏8级地震的最大振幅是里氏5级地震最大振幅的1000倍．
\end{ocrdraft}
\end{bankitem}
